% ============================================================================
% BTGenerator — Architecture d'inférence
% Document de référence pour le déploiement et le service d'inférence
% ============================================================================
\documentclass[a4paper,11pt]{article}

% --- Encodage et langue ---
\usepackage[utf8]{inputenc}
\usepackage[T1]{fontenc}
\usepackage[french]{babel}

% --- Mise en page ---
\usepackage[margin=2.5cm]{geometry}
\usepackage{parskip}            % Espacement entre paragraphes au lieu d'indentation
\usepackage{setspace}
\onehalfspacing

% --- Typographie et polices ---
\usepackage{lmodern}
\usepackage{microtype}

% --- Tableaux ---
\usepackage{booktabs}
\usepackage{array}
\usepackage{longtable}
\usepackage{tabularx}
\usepackage{multirow}
\usepackage{float}

% --- Couleurs et liens ---
\usepackage[dvipsnames]{xcolor}
\usepackage{hyperref}
\hypersetup{
    colorlinks=true,
    linkcolor=NavyBlue,
    citecolor=NavyBlue,
    urlcolor=NavyBlue,
    pdftitle={BTGenerator -- Architecture d'inférence},
    pdfauthor={Benjamin Lepourtois},
}

% --- Code source ---
\usepackage{listings}
\lstset{
    basicstyle=\ttfamily\small,
    breaklines=true,
    frame=single,
    backgroundcolor=\color{gray!8},
    keywordstyle=\color{NavyBlue}\bfseries,
    commentstyle=\color{OliveGreen},
    stringstyle=\color{BrickRed},
    showstringspaces=false,
    tabsize=2,
    xleftmargin=1em,
    framexleftmargin=0.5em,
}

\lstdefinestyle{bash}{
    language=bash,
    morekeywords={ssh,sbatch,curl,sudo,nmcli,uvicorn,pip,scp},
}

\lstdefinestyle{python}{
    language=Python,
    morekeywords={self,def,class,return,from,import,None,True,False},
}

\lstdefinestyle{json}{
    basicstyle=\ttfamily\small,
    string=[s]{"}{"},
    stringstyle=\color{BrickRed},
    comment=[l]{//},
    commentstyle=\color{OliveGreen},
}

% --- Graphiques et diagrammes ---
\usepackage{tikz}
\usetikzlibrary{shapes.geometric, arrows.meta, positioning, calc, fit, backgrounds}

% --- Listes ---
\usepackage{enumitem}

% --- En-tête / pied de page ---
\usepackage{fancyhdr}
\pagestyle{fancy}
\fancyhf{}
\fancyhead[L]{\small\textsc{BTGenerator}}
\fancyhead[R]{\small Architecture d'inférence}
\fancyfoot[C]{\thepage}

% --- Commandes personnalisées ---
\newcommand{\file}[1]{\texttt{#1}}
\newcommand{\skill}[1]{\texttt{<#1/>}}
\newcommand{\tech}[1]{\textsl{#1}}
\newcommand{\apiep}[1]{\texttt{#1}}

% ============================================================================
\begin{document}

% --- Page de titre ---
\begin{titlepage}
\centering
\vspace*{3cm}
{\Huge\bfseries BTGenerator\par}
\vspace{0.5cm}
{\LARGE Architecture d'inférence\par}
\vspace{1.5cm}
{\large\itshape Inférence GPU sur le cluster Télécom Paris\\
avec fallback CPU sur Raspberry Pi~5\\
pour la génération de Behavior Trees NAV4RAIL\par}
\vspace{2cm}
{\large Benjamin Lepourtois --- Télécom Paris\par}
\vspace{1cm}
{\large\today\par}
\vfill
{\small Ce document décrit l'architecture de déploiement du service
d'inférence~: backend GPU primaire sur le cluster SLURM de
Télécom Paris ($\sim$10\,s par requête), fallback CPU sur Raspberry Pi~5
($\sim$6\,min), connectivité réseau dual-homing, et décodage contraint
par grammaire GBNF.}
\end{titlepage}

% --- Table des matières ---
\tableofcontents
\newpage

% ============================================================================
\section{Introduction}
% ============================================================================

Le projet \textbf{BTGenerator} génère des \textbf{Behavior Trees} (BT) au
format XML BehaviorTree.CPP~v4 à partir de descriptions de missions en
langage naturel. Le modèle sous-jacent est un \textbf{Mistral-7B} fine-tuné
par LoRA sur un dataset NAV4RAIL de 2\,000~exemples, puis quantifié au
format \textbf{GGUF Q4\_K\_M} ($\sim$4.1\,Go).

L'inférence repose sur \textbf{llama-cpp-python} et s'exécute sur deux
backends selon la disponibilité~:

\begin{itemize}
    \item \textbf{Backend primaire~:} GPU Tesla P100 (16\,Go VRAM) sur le
    cluster SLURM de Télécom Paris --- $\sim$10--12\,s par génération.
    \item \textbf{Backend de secours~:} CPU ARM Cortex-A76 sur Raspberry
    Pi~5 --- $\sim$5--6\,min par génération.
\end{itemize}

La webapp FastAPI hébergée sur la RPi5 orchestre transparentement le choix
du backend et expose une interface web unique à l'utilisateur.

% ============================================================================
\section{Vue d'ensemble de l'architecture}
\label{sec:overview}
% ============================================================================

\begin{figure}[H]
\centering
\begin{tikzpicture}[
    node distance=1.2cm and 2cm,
    box/.style={rectangle, draw, rounded corners=3pt, minimum width=2.2cm,
                minimum height=1cm, align=center, font=\footnotesize},
    cluster/.style={box, fill=red!8, draw=BrickRed},
    rpi/.style={box, fill=NavyBlue!8, draw=NavyBlue},
    lan/.style={box, fill=OliveGreen!8, draw=OliveGreen},
    arr/.style={-{Stealth[length=5pt]}, thick},
    darr/.style={-{Stealth[length=5pt]}, thick, dashed},
]
    % --- LAN ---
    \node[lan] (browser) {Navigateur\\(LAN WiFi)};

    % --- RPi5 ---
    \node[rpi, right=1.5cm of browser] (app) {\file{app.py}\\FastAPI :5000};
    \node[rpi, below=0.8cm of app] (inf) {\file{inference.py}\\CPU fallback};

    % --- Cluster ---
    \node[cluster, right=1.5cm of app] (gw) {\texttt{gpu-gw}\\(gateway)};
    \node[cluster, right=1.2cm of gw] (gpu) {Nœud GPU\\P100 / 3090\\\file{serve.py} :8080};

    % --- Cadres ---
    \begin{scope}[on background layer]
        \node[draw=NavyBlue, dashed, rounded corners, inner sep=8pt,
              fit=(app)(inf), label={[NavyBlue, font=\small]above:Raspberry Pi 5}] {};
        \node[draw=BrickRed, dashed, rounded corners, inner sep=8pt,
              fit=(gw)(gpu), label={[BrickRed, font=\small]above:Cluster Télécom Paris}] {};
    \end{scope}

    % --- Flèches ---
    \draw[arr] (browser) -- node[above, font=\scriptsize]{HTTP} (app);
    \draw[arr] (app) -- node[above, font=\scriptsize]{SSH tunnel}
                        node[below, font=\scriptsize]{VPN / eth0} (gw);
    \draw[arr] (gw) -- node[above, font=\scriptsize]{Reverse}
                       node[below, font=\scriptsize]{tunnel} (gpu);
    \draw[darr] (app) -- node[right, font=\scriptsize]{Fallback} (inf);
\end{tikzpicture}
\caption{Architecture d'inférence --- vue d'ensemble.
         Trait plein~: chemin primaire (GPU, $\sim$10\,s).
         Trait pointillé~: fallback CPU ($\sim$6\,min).}
\label{fig:overview}
\end{figure}

\begin{table}[H]
\centering
\caption{Comparaison des deux backends d'inférence}
\label{tab:backends}
\begin{tabularx}{\textwidth}{@{}l l l c X@{}}
\toprule
\textbf{Backend} & \textbf{Matériel} & \textbf{Temps} & \textbf{Rôle} & \textbf{Conditions} \\
\midrule
Cluster Télécom & Tesla P100 16\,Go & $\sim$10--12\,s & Primaire & VPN actif + tunnel SSH + job SLURM \\
RPi5 locale & ARM Cortex-A76 & $\sim$5--6\,min & Fallback & Modèle GGUF présent localement \\
\bottomrule
\end{tabularx}
\end{table}

% ============================================================================
\section{Composants}
\label{sec:components}
% ============================================================================

% ----------------------------------------------------------------------------
\subsection{Webapp FastAPI (\file{app.py} --- RPi5)}
\label{sec:webapp}
% ----------------------------------------------------------------------------

Point d'entrée unique pour l'utilisateur. La webapp est hébergée sur la
Raspberry Pi~5, accessible via WiFi sur le LAN local (port~5000).

\subsubsection{Logique de sélection du backend}

À chaque requête \apiep{POST /api/generate}, la webapp tente d'abord
de joindre le cluster~:

\begin{figure}[H]
\centering
\begin{tikzpicture}[
    node distance=0.8cm and 1.2cm,
    decision/.style={diamond, draw, aspect=2.5, inner sep=1pt,
                     align=center, font=\footnotesize},
    process/.style={rectangle, draw, rounded corners=3pt, minimum width=2.5cm,
                    align=center, font=\footnotesize},
    io/.style={rectangle, draw, rounded corners=3pt, minimum width=2.5cm,
               align=center, font=\footnotesize, fill=gray!10},
    gpu/.style={process, fill=OliveGreen!15},
    cpu/.style={process, fill=BrickRed!15},
    arr/.style={-{Stealth[length=5pt]}, thick},
]
    \node[io] (req) {\apiep{POST /api/generate}};
    \node[decision, below=0.8cm of req] (check)
        {Cluster accessible~?\\{\scriptsize\apiep{GET localhost:8080/health}}};
    \node[gpu, below left=1.2cm and 1.2cm of check] (cluster)
        {\apiep{POST} cluster\\{\scriptsize\texttt{/generate}}};
    \node[decision, below right=1.2cm and 1.2cm of check] (loaded)
        {Modèle local\\chargé~?};
    \node[cpu, right=1.2cm of loaded] (load)
        {Chargement GGUF\\{\scriptsize$\sim$2\,min, 5.5\,Go}};
    \node[cpu, below=1.2cm of loaded] (local) {Inférence CPU\\{\scriptsize$\sim$5--6\,min}};
    \node[gpu, below=0.8cm of cluster] (fast) {Réponse\\{\scriptsize$\sim$10--12\,s}};
    \node[io, below=4.5cm of check] (resp) {Retour client\\{\scriptsize(xml, score, errors)}};

    \draw[arr] (req) -- (check);
    \draw[arr] (check) -| node[above left, font=\footnotesize]{oui} (cluster);
    \draw[arr] (check) -| node[above right, font=\footnotesize]{non} (loaded);
    \draw[arr] (loaded) -- node[left, font=\footnotesize]{oui} (local);
    \draw[arr] (loaded) -- node[above, font=\footnotesize]{non} (load);
    \draw[arr] (load) |- (local);
    \draw[arr] (cluster) -- (fast);
    \draw[arr] (fast) |- (resp);
    \draw[arr] (local) |- (resp);
\end{tikzpicture}
\caption{Logique de sélection du backend à chaque requête de génération.}
\label{fig:backend-selection}
\end{figure}

\subsubsection{Endpoints de l'API}

\begin{table}[H]
\centering
\caption{Endpoints exposés par la webapp}
\label{tab:endpoints}
\begin{tabularx}{\textwidth}{@{}l l X@{}}
\toprule
\textbf{Route} & \textbf{Méthode} & \textbf{Description} \\
\midrule
\apiep{/}              & GET  & Interface web (template Jinja2) \\
\apiep{/api/status}    & GET  & Backend actif (\texttt{cluster} ou \texttt{local}) + état du modèle \\
\apiep{/api/generate}  & POST & Génère un BT XML (routing automatique cluster/local) \\
\apiep{/api/validate}  & POST & Valide un XML BT existant (sans inférence) \\
\apiep{/api/examples}  & GET  & Liste de missions d'exemple \\
\apiep{/api/history/stream} & GET & SSE temps réel (statut de génération, historique) \\
\bottomrule
\end{tabularx}
\end{table}

% ----------------------------------------------------------------------------
\subsection{Serveur d'inférence GPU (\file{serve.py} --- Cluster)}
\label{sec:serve}
% ----------------------------------------------------------------------------

Service FastAPI minimal déployé via un job SLURM sur un nœud GPU du
cluster de Télécom Paris.

\subsubsection{Pipeline du job SLURM}

Le fichier \file{job\_serve.sh} orchestre le déploiement complet~:

\begin{enumerate}
    \item \textbf{Chargement des modules}~: CUDA~12.4, GCC~11.5, Python~3.11.
    \item \textbf{Activation du venv}~: environnement contenant
    \texttt{llama-cpp-python} compilé avec support CUDA (GPU offload).
    \item \textbf{Reverse tunnel SSH}~: \texttt{ssh -R 8080:localhost:8080
    gpu-gw} --- rend le port~8080 du nœud GPU accessible depuis la passerelle.
    \item \textbf{Lancement du serveur}~: \texttt{uvicorn serve:app} sur le
    port~8080.
\end{enumerate}

\subsubsection{Pipeline d'inférence}

\begin{figure}[H]
\centering
\begin{tikzpicture}[
    node distance=0.6cm,
    pipestep/.style={rectangle, draw, rounded corners=3pt, minimum width=5cm,
                 minimum height=0.8cm, align=center, font=\small},
    arr/.style={-{Stealth[length=5pt]}, thick},
]
    \node[pipestep] (prompt)  {Construction du prompt \texttt{[INST]...[/INST]}};
    \node[pipestep, below=of prompt] (grammar) {Création d'un objet \texttt{LlamaGrammar} \textbf{frais}};
    \node[pipestep, below=of grammar] (llm) {\texttt{llm()} avec grammaire GBNF + GPU offload};
    \node[pipestep, below=of llm] (extract) {\texttt{extract\_xml()} --- extraction \texttt{<root>...</root>}};
    \node[pipestep, below=of extract] (validate) {\texttt{validate\_bt()} --- L1 / L2 / L3};
    \node[pipestep, below=of validate, fill=gray!10] (resp) {Réponse JSON \{xml, score, errors, \ldots\}};

    \draw[arr] (prompt) -- (grammar);
    \draw[arr] (grammar) -- (llm);
    \draw[arr] (llm) -- (extract);
    \draw[arr] (extract) -- (validate);
    \draw[arr] (validate) -- (resp);
\end{tikzpicture}
\caption{Pipeline d'inférence de \file{serve.py} sur le nœud GPU.}
\label{fig:serve-pipeline}
\end{figure}

\subsubsection{Bug critique --- Grammaire GBNF avec llama-cpp-python v0.3.x}

Avec \texttt{llama-cpp-python}~$\geq$~0.3.0, l'objet \texttt{LlamaGrammar}
encapsule un \texttt{llama\_sampler} dont l'état interne n'est \textbf{pas
réinitialisé} entre les appels. Réutiliser le même objet provoque une
\textbf{ignorance silencieuse} de la grammaire --- le modèle génère alors
des noms de skills hallucinés (ex.~: <<~PassMisson~>> au lieu de
<<~PassMission~>>).

\textbf{Correction~:} créer un objet \texttt{LlamaGrammar} frais par
requête~:

\begin{lstlisting}[style=python, caption={Correction du bug grammar v0.3.x}]
# Correct (v0.3.x) : objet frais par requete
def _make_grammar():
    return LlamaGrammar.from_string(GRAMMAR_STR)

output = llm(prompt, grammar=_make_grammar())
\end{lstlisting}

\begin{lstlisting}[style=python, caption={Bug silencieux --- ne pas faire}]
# Bug : grammaire ignoree apres le premier appel
grammar = LlamaGrammar.from_string(GRAMMAR_STR)
output = llm(prompt, grammar=grammar)  # Ignore silencieusement
\end{lstlisting}

Le serveur inclut un \textbf{self-test au démarrage}~: il génère un court
BT avec grammaire et vérifie que tous les tags produits appartiennent à la
whitelist des 27~skills.

% ----------------------------------------------------------------------------
\subsection{Moteur d'inférence local (\file{inference.py} --- RPi5)}
\label{sec:local}
% ----------------------------------------------------------------------------

Classe \texttt{Nav4RailGenerator} chargeant le modèle GGUF directement en
mémoire CPU sur la RPi5. Activé uniquement si le cluster est inaccessible.

\begin{table}[H]
\centering
\caption{Caractéristiques de l'inférence locale (RPi5)}
\label{tab:local-specs}
\begin{tabularx}{\textwidth}{@{}l X@{}}
\toprule
\textbf{Paramètre} & \textbf{Valeur} \\
\midrule
Modèle         & \file{nav4rail-mistral-7b-q4\_k\_m.gguf} ($\sim$4.1\,Go, Q4\_K\_M) \\
Contexte       & 2\,048 tokens \\
Threads        & Auto-détecté (4~cœurs ARM Cortex-A76) \\
RAM utilisée   & $\sim$5.5\,Go \\
Chargement     & $\sim$2\,min \\
Inférence      & $\sim$5--6\,min par mission \\
\bottomrule
\end{tabularx}
\end{table}

% ============================================================================
\section{Réseau et connectivité}
\label{sec:network}
% ============================================================================

% ----------------------------------------------------------------------------
\subsection{Dual-homing de la Raspberry Pi~5}
\label{sec:dualhoming}
% ----------------------------------------------------------------------------

La RPi5 utilise deux interfaces réseau pour séparer le trafic~:

\begin{figure}[H]
\centering
\begin{tikzpicture}[
    node distance=0.8cm and 1.5cm,
    box/.style={rectangle, draw, rounded corners=3pt, minimum width=2cm,
                minimum height=0.8cm, align=center, font=\footnotesize},
    iface/.style={box, fill=NavyBlue!10, draw=NavyBlue},
    ext/.style={box, fill=gray!10, draw=gray!50},
    arr/.style={-{Stealth[length=5pt]}, thick},
    darr/.style={{Stealth[length=5pt]}-{Stealth[length=5pt]}, thick},
]
    % RPi5
    \node[box, fill=NavyBlue!5] (app) {\file{app.py} :5000};
    \node[iface, below left=1cm and 0.5cm of app] (wlan) {\texttt{wlan0}\\WiFi};
    \node[iface, below right=1cm and 0.5cm of app] (eth) {\texttt{eth0}\\Ethernet};
    \node[iface, below=0.8cm of eth] (tun) {\texttt{tun0}\\VPN};

    % Externes
    \node[ext, left=2cm of wlan] (pc) {Navigateur\\192.168.x.x};
    \node[ext, right=2cm of tun] (gw) {\texttt{gpu-gw}\\137.194.132.200};

    % Cadre RPi5
    \begin{scope}[on background layer]
        \node[draw=NavyBlue, dashed, rounded corners, inner sep=10pt,
              fit=(app)(wlan)(eth)(tun),
              label={[NavyBlue, font=\small]above:Raspberry Pi 5}] {};
    \end{scope}

    % Flèches
    \draw[darr] (pc) -- node[above, font=\scriptsize]{HTTP} (wlan);
    \draw[arr] (wlan) -- (app);
    \draw[arr] (eth) -- (tun);
    \draw[darr] (tun) -- node[above, font=\scriptsize]{OpenVPN} (gw);
    \draw[darr, dashed] (app) to[bend right=30]
        node[left, font=\scriptsize, xshift=-3pt]{SSH :8080} (tun);
\end{tikzpicture}
\caption{Dual-homing~: WiFi pour la webapp, Ethernet + VPN pour le cluster.}
\label{fig:dualhoming}
\end{figure}

\begin{table}[H]
\centering
\caption{Répartition des interfaces réseau}
\label{tab:interfaces}
\begin{tabularx}{\textwidth}{@{}l l X@{}}
\toprule
\textbf{Interface} & \textbf{Réseau} & \textbf{Usage} \\
\midrule
\texttt{wlan0} (WiFi)     & LAN local (192.168.x.x) & Servir la webapp aux navigateurs \\
\texttt{eth0} (Ethernet)  & VPN Télécom Paris (\texttt{tun0}) & Accéder au cluster via \texttt{gpu-gw} \\
\bottomrule
\end{tabularx}
\end{table}

Configuration VPN essentielle~: \texttt{ipv4.never-default yes} empêche
le VPN d'écraser la route par défaut (WiFi). La route vers le réseau
Télécom est ajoutée manuellement~:

\begin{lstlisting}[style=bash]
sudo nmcli connection up telecom-paris
sudo ip route add 137.194.0.0/16 dev tun0
\end{lstlisting}

\textbf{Attention~:} cette route est éphémère et perdue au redémarrage du
VPN. Elle peut être automatisée via un script NetworkManager dispatcher
dans \file{/etc/NetworkManager/dispatcher.d/}.

% ----------------------------------------------------------------------------
\subsection{Chaîne de tunnels SSH}
\label{sec:tunnels}
% ----------------------------------------------------------------------------

Le cluster Télécom Paris n'expose pas de ports publics. Les nœuds GPU ne
sont pas accessibles directement en SSH depuis l'extérieur
(\texttt{pam\_slurm\_adopt} bloque les connexions entrantes). L'accès
repose sur un \textbf{double tunnel SSH}~:

\begin{figure}[H]
\centering
\begin{tikzpicture}[
    node distance=1.5cm and 3cm,
    box/.style={rectangle, draw, rounded corners=3pt, minimum width=2.5cm,
                minimum height=1.2cm, align=center, font=\small},
    arr/.style={-{Stealth[length=6pt]}, thick},
    tunnel/.style={arr, NavyBlue},
    data/.style={arr, OliveGreen, dashed},
]
    \node[box] (rpi) {RPi5\\{\footnotesize localhost:8080}};
    \node[box, right=3.5cm of rpi] (gw) {\texttt{gpu-gw}\\{\footnotesize :8080}};
    \node[box, right=3.5cm of gw] (node) {Nœud GPU\\{\footnotesize serve.py :8080}};

    % Tunnels
    \draw[tunnel] (rpi) -- node[above, font=\footnotesize]
        {\texttt{ssh -L 8080:localhost:8080}} (gw);
    \draw[tunnel] (node) -- node[above, font=\footnotesize]
        {\texttt{ssh -R 8080:localhost:8080}} (gw);

    % Labels
    \node[below=0.3cm of rpi, font=\footnotesize\itshape, NavyBlue]
        {Forward tunnel};
    \node[below=0.3cm of node, font=\footnotesize\itshape, NavyBlue]
        {Reverse tunnel};

    % Flux de données
    \draw[data, bend left=25] (rpi) to
        node[below, font=\footnotesize, yshift=-5pt]
        {\texttt{curl localhost:8080/generate} $\longrightarrow$ GPU}
        (node);
\end{tikzpicture}
\caption{Double tunnel SSH~: le flux de données traverse
         \texttt{gpu-gw} de manière transparente.}
\label{fig:tunnels}
\end{figure}

\begin{enumerate}
    \item \textbf{Reverse tunnel (depuis le job SLURM)~:} le nœud GPU ouvre
    un tunnel vers \texttt{gpu-gw} au démarrage du job (\texttt{ssh~-R}).
    C'est le seul sens possible car les connexions entrantes sont bloquées
    par \texttt{pam\_slurm\_adopt}.

    \item \textbf{Forward tunnel (depuis la RPi5)~:} la RPi5 mappe son
    \texttt{localhost:8080} vers le port~8080 de \texttt{gpu-gw}
    (\texttt{ssh~-L}), qui est relié au nœud GPU par le reverse tunnel.
\end{enumerate}

\textbf{Résultat~:} \texttt{curl http://localhost:8080/generate} sur la
RPi5 atteint directement le GPU du cluster.

% ============================================================================
\section{Modèle et décodage contraint}
\label{sec:model}
% ============================================================================

% ----------------------------------------------------------------------------
\subsection{Modèle}
% ----------------------------------------------------------------------------

\begin{table}[H]
\centering
\caption{Caractéristiques du modèle d'inférence}
\label{tab:model}
\begin{tabularx}{\textwidth}{@{}l X@{}}
\toprule
\textbf{Propriété} & \textbf{Valeur} \\
\midrule
Base         & Mistral-7B-Instruct-v0.2 \\
Fine-tuning  & LoRA ($r=16$, $\alpha=32$) sur dataset NAV4RAIL (2\,000 exemples) \\
Format       & GGUF Q4\_K\_M ($\sim$4.1\,Go, 4.83~BPW) \\
Backend      & llama-cpp-python (CPU sur RPi5, CUDA sur cluster) \\
Contexte     & 2\,048 tokens \\
\bottomrule
\end{tabularx}
\end{table}

% ----------------------------------------------------------------------------
\subsection{Grammaire GBNF}
\label{sec:gbnf}
% ----------------------------------------------------------------------------

La grammaire GBNF (\file{nav4rail\_grammar.py}) contraint le décodage
\textbf{token par token}. À chaque étape de génération, seuls les tokens
compatibles avec la grammaire sont autorisés. Elle garantit~:

\begin{itemize}
    \item \textbf{Structure XML valide~:} enveloppe
    \texttt{<root BTCPP\_format="4">} contenant un \texttt{<Behavior\-Tree ID="MainTree">}.
    \item \textbf{Nœuds de contrôle~:} uniquement \texttt{Sequence} et
    \texttt{Fallback}, avec imbrication libre.
    \item \textbf{Nœuds feuilles~:} uniquement les \textbf{27~skills
    NAV4RAIL} (4~familles~: Preparation, Motion, Inspection, Simulation).
    \item \textbf{Attributs~:} \texttt{name} en \texttt{snake\_case}.
    \item \textbf{Zéro hallucination} de nom de skill.
\end{itemize}

Extrait de la grammaire~:

\begin{lstlisting}[caption={Règle GBNF des skills autorisés (extrait)}]
skilltag ::= "LoadMission" | "MissionStructureValid"
           | "UpdateCurrentGeneratedActivity"
           | "ProjectPointOnNetwork" | "CreatePath"
           | "AgregatePath" | ... | "SimulationStarted"
\end{lstlisting}

% ----------------------------------------------------------------------------
\subsection{Validation multi-niveaux (\file{validate\_bt.py})}
\label{sec:validation}
% ----------------------------------------------------------------------------

Appliquée systématiquement après chaque génération, quel que soit le
backend~:

\begin{figure}[H]
\centering
\begin{tikzpicture}[
    node distance=0.6cm and 1cm,
    level/.style={rectangle, draw, rounded corners=3pt, minimum width=3.2cm,
                  minimum height=1.3cm, align=center, font=\footnotesize},
    result/.style={rectangle, draw, rounded corners=3pt, minimum width=1.8cm,
                   minimum height=0.7cm, align=center, font=\footnotesize\bfseries},
    arr/.style={-{Stealth[length=5pt]}, thick},
]
    \node[level, fill=NavyBlue!10] (l1)
        {\textbf{L1 --- Syntaxique}\\XML bien formé\\Tags autorisés, format v4};
    \node[level, fill=OliveGreen!10, right=1.5cm of l1] (l2)
        {\textbf{L2 --- Structurel}\\Nœuds non vides\\Fallback $\geq$ 2 branches\\Profondeur $\leq$ 10};
    \node[level, fill=Goldenrod!15, right=1.5cm of l2] (l3)
        {\textbf{L3 --- Sémantique}\\Ordre des skills\\Patterns métier};

    \node[result, fill=OliveGreen!20, below right=1.2cm and -0.5cm of l3] (ok)
        {Score = 1.0};
    \node[result, fill=BrickRed!20, below=2cm of l1] (ko) {Score = 0.0};

    \draw[arr] (l1) -- node[above, font=\footnotesize]{OK} (l2);
    \draw[arr] (l2) -- node[above, font=\footnotesize]{OK} (l3);
    \draw[arr] (l3) -- node[right, font=\footnotesize]{OK} (ok);
    \draw[arr, BrickRed] (l1) -- node[left, font=\footnotesize]{Erreur} (ko);
    \draw[arr, BrickRed] (l2) |- (ko);
    \draw[arr, BrickRed] (l3) |- node[above, font=\footnotesize, near start]{Erreur} (ko);
\end{tikzpicture}
\caption{Pipeline de validation à trois niveaux. Pénalité de $-0.1$ par
         warning sémantique.}
\label{fig:validation}
\end{figure}

% ============================================================================
\section{Contrat API du service d'inférence}
\label{sec:api-contract}
% ============================================================================

Le service GPU (\file{serve.py}) et le moteur local (\file{inference.py})
partagent le même format de réponse, permettant au frontend de traiter les
deux backends de manière transparente~:

\begin{lstlisting}[style=json, caption={Format de réponse commun aux deux backends}]
{
  "xml": "<root BTCPP_format=\"4\">...</root>",
  "valid": true,
  "score": 1.0,
  "errors": [],
  "warnings": [],
  "summary": "OK (score=1.0) -- L1+L2+L3 passes",
  "generation_time_s": 11.5
}
\end{lstlisting}

% ============================================================================
\section{Déploiement}
\label{sec:deployment}
% ============================================================================

\subsection{Lancer le service GPU sur le cluster}

\begin{lstlisting}[style=bash, caption={Soumission du job SLURM depuis la RPi5}]
# Depuis la RPi5 (VPN actif)
ssh gpu
cd ~/nav4rail_serve
sbatch job_serve.sh

# Surveiller le demarrage
tail -f nav4rail_serve_<jobid>.out
# Attendre : "[serve] Grammar self-test PASSED"
\end{lstlisting}

\subsection{Ouvrir le tunnel depuis la RPi5}

\begin{lstlisting}[style=bash]
ssh -fN -L 8080:localhost:8080 gpu
\end{lstlisting}

\subsection{Lancer la webapp}

\begin{lstlisting}[style=bash]
cd ~/nav4rail_webapp
source .venv/bin/activate
uvicorn app:app --host 0.0.0.0 --port 5000
\end{lstlisting}

\subsection{Vérification end-to-end}

\begin{lstlisting}[style=bash, caption={Tests de bout en bout}]
# Health du cluster
curl -s localhost:8080/health
# -> {"status":"ok","gpu":true,"model":"nav4rail-mistral-7b-q4_k_m.gguf"}

# Inference via la webapp
curl -s -X POST localhost:5000/api/generate \
  -H 'Content-Type: application/json' \
  -d '{"mission":"Navigue au km 42 depuis le km 10"}'
# -> {"xml":"<root ...>...</root>","score":1.0,"generation_time_s":11.5}
\end{lstlisting}

% ============================================================================
\section{Limitations et points d'attention}
\label{sec:limitations}
% ============================================================================

\begin{table}[H]
\centering
\caption{Limitations connues et mitigations}
\label{tab:limitations}
\begin{tabularx}{\textwidth}{@{}l X@{}}
\toprule
\textbf{Sujet} & \textbf{Détail} \\
\midrule
Job SLURM temporaire &
    Le service GPU est un job de 2\,h max. Il faut le relancer
    périodiquement via \texttt{sbatch}. \\
Route VPN éphémère &
    \texttt{ip route add 137.194.0.0/16 dev tun0} est perdu au
    redémarrage du VPN. Automatiser via un dispatcher NetworkManager. \\
Tunnel SSH fragile &
    Si le tunnel tombe, la webapp bascule silencieusement sur le CPU
    local ($\sim$6\,min par requête). \\
Modèle GGUF unique &
    Le même fichier \texttt{.gguf} doit être présent sur le cluster ET
    la RPi5 pour des résultats cohérents. \\
Grammar v0.3.x &
    Créer un objet \texttt{LlamaGrammar} frais par requête (cf.\
    section~\ref{sec:serve}). \\
\bottomrule
\end{tabularx}
\end{table}

\end{document}
